\documentclass[tikz,border=6pt]{standalone}
\usepackage[no-math]{fontspec}
\usepackage[brazil]{babel}
\usepackage{amsmath,amssymb}
\IfFileExists{C:/Windows/Fonts/Lato_Regular.ttf}{%
  \setmainfont{Lato}[
    Path = C:/Windows/Fonts/,
    UprightFont = Lato_Regular.ttf,
    BoldFont = Lato_Bold.ttf,
    ItalicFont = Lato_Italic.ttf,
    BoldItalicFont = Lato_BoldItalic.ttf
  ]
  \setsansfont{Lato}[
    Path = C:/Windows/Fonts/,
    UprightFont = Lato_Regular.ttf,
    BoldFont = Lato_Bold.ttf,
    ItalicFont = Lato_Italic.ttf,
    BoldItalicFont = Lato_BoldItalic.ttf
  ]
}{%
  \RequirePackage[default]{lato}
}
\renewcommand{\familydefault}{\sfdefault}

\usepackage{xcolor}
\usepackage{tikz}
\usepackage{pgfplots}
\pgfplotsset{compat=1.18}
\usetikzlibrary{arrows.meta,positioning,shapes.geometric,fit,backgrounds,calc,decorations.pathreplacing,patterns}

% Paleta institucional UEFS / CHF662
\definecolor{uefsNavy}{HTML}{1B2A4A}
\definecolor{uefsTeal}{HTML}{0E7C7B}
\definecolor{uefsForest}{HTML}{2D6A4F}
\definecolor{uefsGold}{HTML}{C4983B}
\definecolor{uefsGoldDark}{HTML}{8F6B1E}
\definecolor{uefsAlert}{HTML}{B02A37}
\definecolor{uefsOrange}{HTML}{D95F02}
\definecolor{uefsInk}{HTML}{22252A}
\definecolor{uefsMuted}{HTML}{5C6470}
\definecolor{uefsWater}{HTML}{2E86AB}
\definecolor{uefsSand}{HTML}{D9B382}
\definecolor{bgNavy}{HTML}{F2F5F9}
\definecolor{bgTeal}{HTML}{E9F5F4}
\definecolor{bgForest}{HTML}{EEF5F0}
\definecolor{bgGold}{HTML}{FCF7EB}
\definecolor{bgAlert}{HTML}{FDF2F2}
\definecolor{bgGray}{HTML}{F8FAFC}
\definecolor{lineGray}{HTML}{CFD8DC}

\tikzset{
  >=Stealth,
  aaCard/.style={draw=#1, fill=white, rounded corners=4pt, line width=1pt, inner sep=6pt, align=center, font=\small},
  aaBox/.style={draw=#1, rounded corners=4pt, line width=1.1pt, inner sep=6pt, align=center, font=\footnotesize\sffamily, text=uefsInk},
  aaTitle/.style={font=\normalsize\bfseries\sffamily, text=uefsNavy, align=center},
  aaSub/.style={font=\footnotesize\sffamily, text=uefsMuted, align=center},
  aaArrow/.style={->, >=Stealth, line width=1.4pt},
  aaDArrow/.style={<->, >=Stealth, line width=1.4pt}
}


\begin{document}
\begin{tikzpicture}[font=\sffamily]

% =========================================================
% PAINEL (a) Cadeia Causal na Análise de Impactos (Sánchez, 2013)
% =========================================================
\node[aaTitle] at (0, 3.25) {(a) Cadeia Causal na Análise de Impactos (Sánchez, 2013)};

\node[aaBox=uefsNavy, fill=bgNavy, text width=3.35cm, minimum height=2.75cm, inner sep=7pt] (ativ) at (-4.1, 1.05) {
  \textbf{\color{uefsNavy}ATIVIDADE}\\[-1pt]
  \textbf{\color{uefsNavy}HUMANA}\\[1pt]
  \textbf{(Causa)}\\[4pt]
  \scriptsize Ex.: abertura de estrada em encosta
};

\node[aaBox=uefsGold, fill=bgGold, text width=3.35cm, minimum height=2.75cm, inner sep=7pt] (asp) at (0, 1.05) {
  \textbf{\color{uefsGoldDark}ASPECTO}\\[-1pt]
  \textbf{\color{uefsGoldDark}AMBIENTAL}\\[1pt]
  \textbf{(Mecanismo)}\\[4pt]
  \scriptsize Ex.: exposição do solo à chuva
};

\node[aaBox=uefsAlert, fill=bgAlert, text width=3.35cm, minimum height=2.75cm, inner sep=7pt] (imp) at (4.1, 1.05) {
  \textbf{\color{uefsAlert}IMPACTO}\\[-1pt]
  \textbf{\color{uefsAlert}AMBIENTAL}\\[1pt]
  \textbf{(Efeito)}\\[4pt]
  \scriptsize Ex.: erosão linear e assoreamento
};

\draw[aaArrow, uefsNavy, shorten >=3pt, shorten <=3pt] (ativ.east) -- (asp.west);
\draw[aaArrow, uefsAlert, shorten >=3pt, shorten <=3pt] (asp.east) -- (imp.west);

% Quadro inferior amplo contendo a distinção de Sánchez (2013)
\node[aaBox=uefsTeal, fill=bgTeal, text width=11.4cm, minimum height=1.65cm, inner sep=9pt] (dist) at (0, -1.75) {
  \footnotesize\textbf{\color{uefsNavy}Distinção essencial (Sánchez, 2013):} o \textbf{ASPECTO} é o elemento da atividade que interage com o ambiente; o \textbf{IMPACTO} é a alteração resultante na qualidade ambiental (pode ser adverso ou benéfico).
};

% =========================================================
% PAINEL (b) Matriz Conceitual de Risco (R = P x V)
% =========================================================
\node[aaTitle] at (12.2, 3.25) {(b) Matriz Conceitual de Risco ($R = P \times V$)};

\begin{scope}[shift={(9.3, -2.3)}]
  % Cores graduais por nível de risco
  \definecolor{rBaixo1}{HTML}{D8F3DC}
  \definecolor{rBaixo2}{HTML}{B7E4C7}
  \definecolor{rBaixo3}{HTML}{FFF3B0}
  \definecolor{rMedio}{HTML}{FFE066}
  \definecolor{rAlto}{HTML}{FDAE61}
  \definecolor{rCrit1}{HTML}{F46D43}
  \definecolor{rCrit2}{HTML}{D73027}

  % Tamanho de cada célula: 1.45 x 1.2
  \def\cw{1.45}
  \def\ch{1.2}

  % Linha P=1 (y=0)
  \filldraw[fill=rBaixo1, draw=white, line width=1.2pt] (0*\cw, 0*\ch) rectangle (1*\cw, 1*\ch) node[midway, align=center, font=\scriptsize\bfseries, text=uefsInk] {1\\(Baixo)};
  \filldraw[fill=rBaixo2, draw=white, line width=1.2pt] (1*\cw, 0*\ch) rectangle (2*\cw, 1*\ch) node[midway, align=center, font=\scriptsize\bfseries, text=uefsInk] {2\\(Baixo)};
  \filldraw[fill=rBaixo3, draw=white, line width=1.2pt] (2*\cw, 0*\ch) rectangle (3*\cw, 1*\ch) node[midway, align=center, font=\scriptsize\bfseries, text=uefsInk] {3\\(Baixo)};
  \filldraw[fill=rMedio,  draw=white, line width=1.2pt] (3*\cw, 0*\ch) rectangle (4*\cw, 1*\ch) node[midway, align=center, font=\scriptsize\bfseries, text=uefsInk] {4\\(Médio)};

  % Linha P=2 (y=1)
  \filldraw[fill=rBaixo2, draw=white, line width=1.2pt] (0*\cw, 1*\ch) rectangle (1*\cw, 2*\ch) node[midway, align=center, font=\scriptsize\bfseries, text=uefsInk] {2\\(Baixo)};
  \filldraw[fill=rMedio,  draw=white, line width=1.2pt] (1*\cw, 1*\ch) rectangle (2*\cw, 2*\ch) node[midway, align=center, font=\scriptsize\bfseries, text=uefsInk] {4\\(Médio)};
  \filldraw[fill=rAlto!65!rMedio, draw=white, line width=1.2pt] (2*\cw, 1*\ch) rectangle (3*\cw, 2*\ch) node[midway, align=center, font=\scriptsize\bfseries, text=uefsInk] {6\\(Médio)};
  \filldraw[fill=rAlto,   draw=white, line width=1.2pt] (3*\cw, 1*\ch) rectangle (4*\cw, 2*\ch) node[midway, align=center, font=\scriptsize\bfseries, text=uefsInk] {8\\(Alto)};

  % Linha P=3 (y=2)
  \filldraw[fill=rBaixo3, draw=white, line width=1.2pt] (0*\cw, 2*\ch) rectangle (1*\cw, 3*\ch) node[midway, align=center, font=\scriptsize\bfseries, text=uefsInk] {3\\(Baixo)};
  \filldraw[fill=rAlto!65!rMedio, draw=white, line width=1.2pt] (1*\cw, 2*\ch) rectangle (2*\cw, 3*\ch) node[midway, align=center, font=\scriptsize\bfseries, text=uefsInk] {6\\(Médio)};
  \filldraw[fill=rCrit1!75!white, draw=white, line width=1.2pt] (2*\cw, 2*\ch) rectangle (3*\cw, 3*\ch) node[midway, align=center, font=\scriptsize\bfseries, text=uefsInk] {9\\(Alto)};
  \filldraw[fill=rCrit1,  draw=white, line width=1.2pt] (3*\cw, 2*\ch) rectangle (4*\cw, 3*\ch) node[midway, align=center, font=\scriptsize\bfseries, text=white] {12\\(Crítico)};

  % Linha P=4 (y=3)
  \filldraw[fill=rMedio,  draw=white, line width=1.2pt] (0*\cw, 3*\ch) rectangle (1*\cw, 4*\ch) node[midway, align=center, font=\scriptsize\bfseries, text=uefsInk] {4\\(Médio)};
  \filldraw[fill=rAlto,   draw=white, line width=1.2pt] (1*\cw, 3*\ch) rectangle (2*\cw, 4*\ch) node[midway, align=center, font=\scriptsize\bfseries, text=uefsInk] {8\\(Alto)};
  \filldraw[fill=rCrit1,  draw=white, line width=1.2pt] (2*\cw, 3*\ch) rectangle (3*\cw, 4*\ch) node[midway, align=center, font=\scriptsize\bfseries, text=white] {12\\(Crítico)};
  \filldraw[fill=rCrit2,  draw=white, line width=1.2pt] (3*\cw, 3*\ch) rectangle (4*\cw, 4*\ch) node[midway, align=center, font=\scriptsize\bfseries, text=white] {16\\(Crítico)};

  % Borda externa da matriz
  \draw[lineGray, line width=1pt] (0,0) rectangle (4*\cw, 4*\ch);

  % Ticks e Rótulos do Eixo X (Vulnerabilidade V)
  \foreach \i/\lbl in {0.5/{1\\Baixa}, 1.5/{2\\Média}, 2.5/{3\\Alta}, 3.5/{4\\Muito alta}} {
    \draw[uefsInk, line width=0.8pt] (\i*\cw, 0) -- (\i*\cw, -0.1);
    \node[below=3pt, align=center, font=\scriptsize, text=uefsInk] at (\i*\cw, -0.08) {\lbl};
  }
  \node[below=26pt, font=\footnotesize\bfseries, text=uefsInk] at (2*\cw, -0.05) {Vulnerabilidade ($V$)};

  % Ticks e Rótulos do Eixo Y (Perigo / Ameaça P)
  \foreach \j/\lbl in {0.5/{1 Baixo}, 1.5/{2 Médio}, 2.5/{3 Alto}, 3.5/{4 Crítico}} {
    \draw[uefsInk, line width=0.8pt] (0, \j*\ch) -- (-0.1, \j*\ch);
    \node[left=3pt, align=right, font=\scriptsize, text=uefsInk] at (-0.08, \j*\ch) {\lbl};
  }
  \node[rotate=90, anchor=center, font=\footnotesize\bfseries, text=uefsInk] at (-1.58, 2*\ch) {Perigo / Ameaça ($P$)};
\end{scope}

\end{tikzpicture}
\end{document}
